\documentclass[10pt,a4paper]{amsart}
\usepackage{amsmath,amssymb,amsthm,graphicx}
\usepackage{verbatim}
\newenvironment{reference}{\comment}{\endcomment}
%\newenvironment{reference}{}{}
\newenvironment{slogan}{\comment}{\endcomment}
\newenvironment{history}{\comment}{\endcomment}

% For commutative diagrams we use Xy-pic
\usepackage[all]{xy}

% We use 2cell for 2-commutative diagrams.
\xyoption{2cell}
% Theorem environments.
%
\theoremstyle{plain}
\newtheorem{theorem}[subsection]{Theorem}
\newtheorem{proposition}[subsection]{Proposition}
\newtheorem{lemma}[subsection]{Lemma}

\theoremstyle{definition}
\newtheorem{definition}[subsection]{Definition}
\newtheorem{example}[subsection]{Example}
\newtheorem{exercise}[subsection]{Exercise}
\newtheorem{situation}[subsection]{Situation}

\theoremstyle{remark}
\newtheorem{remark}[subsection]{Remark}
\newtheorem{remarks}[subsection]{Remarks}

\numberwithin{equation}{subsection}

% Macros
%
\def\lim{\mathop{\mathrm{lim}}\nolimits}
\def\colim{\mathop{\mathrm{colim}}\nolimits}
\def\Spec{\mathop{\mathrm{Spec}}}
\def\Hom{\mathop{\mathrm{Hom}}\nolimits}
\def\Ext{\mathop{\mathrm{Ext}}\nolimits}
\def\SheafHom{\mathop{\mathcal{H}\!\mathit{om}}\nolimits}
\def\SheafExt{\mathop{\mathcal{E}\!\mathit{xt}}\nolimits}
\def\Sch{\mathit{Sch}}
\def\Mor{\mathop{\mathrm{Mor}}\nolimits}
\def\Ob{\mathop{\mathrm{Ob}}\nolimits}
\def\Sh{\mathop{\mathit{Sh}}\nolimits}
\def\NL{\mathop{N\!L}\nolimits}
\def\CH{\mathop{\mathrm{CH}}\nolimits}
\def\proetale{{pro\text{-}\acute{e}tale}}
\def\etale{{\acute{e}tale}}
\def\QCoh{\mathit{QCoh}}
\def\Ker{\mathop{\mathrm{Ker}}}
\def\Im{\mathop{\mathrm{Im}}}
\def\Coker{\mathop{\mathrm{Coker}}}
\def\Coim{\mathop{\mathrm{Coim}}}

% Boxtimes
%
\DeclareMathSymbol{\boxtimes}{\mathbin}{AMSa}{"02}

%
% Macros for moduli stacks/spaces
%
\def\QCohstack{\mathcal{QC}\!\mathit{oh}}
\def\Cohstack{\mathcal{C}\!\mathit{oh}}
\def\Spacesstack{\mathcal{S}\!\mathit{paces}}
\def\Quotfunctor{\mathrm{Quot}}
\def\Hilbfunctor{\mathrm{Hilb}}
\def\Curvesstack{\mathcal{C}\!\mathit{urves}}
\def\Polarizedstack{\mathcal{P}\!\mathit{olarized}}
\def\Complexesstack{\mathcal{C}\!\mathit{omplexes}}
% \Pic is the operator that assigns to X its picard group, usage \Pic(X)
% \Picardstack_{X/B} denotes the Picard stack of X over B
% \Picardfunctor_{X/B} denotes the Picard functor of X over B
\def\Pic{\mathop{\mathrm{Pic}}\nolimits}
\def\Picardstack{\mathcal{P}\!\mathit{ic}}
\def\Picardfunctor{\mathrm{Pic}}
\def\Deformationcategory{\mathcal{D}\!\mathit{ef}}
\usepackage{cleveref}
\crefname{originallemma}{original lemma}{original lemmas}
\Crefname{originallemma}{Original lemma}{Original lemmas}
\makeatletter
\newlabel{lemma-finite-type-points-morphism}{{06EI}{source}{}{sourcepointer-1}{}}
\newlabel{lemma-finite-type-points-morphism@cref}{{[originallemma][0][]06EI}{[1][0][]source}{}{}{}}
\makeatother
\crefname{originallemma}{original lemma}{original lemmas}
\Crefname{originallemma}{Original lemma}{Original lemmas}
\makeatletter
\newlabel{lemma-identify-finite-type-points}{{06EH}{source}{}{sourcepointer-2}{}}
\newlabel{lemma-identify-finite-type-points@cref}{{[originallemma][0][]06EH}{[1][0][]source}{}{}{}}
\makeatother
\crefname{originallemma}{original lemma}{original lemmas}
\Crefname{originallemma}{Original lemma}{Original lemmas}
\makeatletter
\newlabel{lemma-immersion-locally-finite-type}{{06ED}{source}{}{sourcepointer-3}{}}
\newlabel{lemma-immersion-locally-finite-type@cref}{{[originallemma][0][]06ED}{[1][0][]source}{}{}{}}
\makeatother
\crefname{originalremark}{original remark}{original remarks}
\Crefname{originalremark}{Original remark}{Original remarks}
\makeatletter
\newlabel{remark-space-structure-locally-closed-subset}{{06EC}{source}{}{sourcepointer-4}{}}
\newlabel{remark-space-structure-locally-closed-subset@cref}{{[originalremark][0][]06EC}{[1][0][]source}{}{}{}}
\makeatother
\title{Every nonempty locally closed subset of an algebraic space meets the finite type points}
\author{The Stacks Project Authors}
\date{stacks-project:spaces-morphisms.tex:local-snapshot-20261012; GNU FDL 1.2 or later}
\begin{document}
\maketitle
\setcounter{subsection}{1}
\noindent\textit{Editorial note.} All delivered mathematical text is the author's own, copied from the pinned Stacks Project chapter file spaces-morphisms.tex; no statement or proof step was rewritten, reordered or removed.\par\medskip
\noindent\textit{Editorial note.} The statement includes a displayed implication and an 'in particular' sentence; both are inside the counted statement span.\par\medskip
\noindent\textit{Editorial note.} The four references of the counted spans are recorded as non-load-bearing navigation pointers; the argument that consumes them stays inside the counted proof span.\par\medskip
\noindent\textit{Editorial note.} Preamble: the three native excerpt spans of sources/stacks-spaces-morphisms/preamble.tex (lines 8-13, 14-18 and 154-217). The chapter file itself carries no preamble text.\par\medskip
\noindent\textit{Allowed standard prerequisites (editorial).}\begin{itemize}
\item The reduced induced subspace structure on a locally closed subset. Remark remark-space-structure-locally-closed-subset, used at proof lines 5144-5146. Produces the algebraic space Z with |Z| = T.
\item Immersions are locally of finite type. Lemma lemma-immersion-locally-finite-type, used at proof lines 5147-5148. Makes Z -> X locally of finite type.
\item Finite type points behave well under locally finite type morphisms. Lemma lemma-finite-type-points-morphism, used at proof lines 5149-5151. Gives that the finite type points of Z lie among those of X.
\item Nonempty affine schemes have closed points and finite type points of a space come from them. Lemma lemma-identify-finite-type-points, used at proof lines 5152-5155. Produces a finite type point of Z.
\end{itemize}\medskip
\ifdefined\hypertarget\hypertarget{sourcepointer-1}{}\fi
\noindent\textit{Source reference.} Original lemma 06EI: focused-seeds/AGP/sources/stacks-snapshot-20261011/tags.txt, source lines 8217--8217. Navigation pointer into the Stacks Project; the standard fact it supplies is applied inside the counted proof span and the argument resting on it stays there.\par\medskip
\ifdefined\hypertarget\hypertarget{sourcepointer-2}{}\fi
\noindent\textit{Source reference.} Original lemma 06EH: focused-seeds/AGP/sources/stacks-snapshot-20261011/tags.txt, source lines 8216--8216. Navigation pointer into the Stacks Project; the standard fact it supplies is applied inside the counted proof span and the argument resting on it stays there.\par\medskip
\ifdefined\hypertarget\hypertarget{sourcepointer-3}{}\fi
\noindent\textit{Source reference.} Original lemma 06ED: focused-seeds/AGP/sources/stacks-snapshot-20261011/tags.txt, source lines 8212--8212. Navigation pointer into the Stacks Project; the standard fact it supplies is applied inside the counted proof span and the argument resting on it stays there.\par\medskip
\ifdefined\hypertarget\hypertarget{sourcepointer-4}{}\fi
\noindent\textit{Source reference.} Original remark 06EC: focused-seeds/AGP/sources/stacks-snapshot-20261011/tags.txt, source lines 8211--8211. Navigation pointer into the Stacks Project; the standard fact it supplies is applied inside the counted proof span and the argument resting on it stays there.\par\medskip
\begin{lemma}
\label{lemma-enough-finite-type-points}
Let $S$ be a scheme. Let $X$ be an algebraic space over $S$.
For any locally closed subset $T \subset |X|$ we have
$$
T \not = \emptyset
\Rightarrow
T \cap X_{\text{ft-pts}} \not = \emptyset.
$$
In particular, for any closed subset $T \subset |X|$ we
see that $T \cap X_{\text{ft-pts}}$ is dense in $T$.
\end{lemma}

\begin{proof}
Let $i : Z \to X$ be the reduced induce subspace structure on $T$, see
Remark \ref{remark-space-structure-locally-closed-subset}.
Any immersion is locally of finite type, see
Lemma \ref{lemma-immersion-locally-finite-type}.
Hence by
Lemma \ref{lemma-finite-type-points-morphism}
we see $Z_{\text{ft-pts}} \subset X_{\text{ft-pts}} \cap T$.
Finally, any nonempty affine scheme $U$ with an \'etale morphism towards
$Z$ has at least one closed point. Hence $Z$ has at least one
finite type point by
Lemma \ref{lemma-identify-finite-type-points}.
The lemma follows.
\end{proof}

\end{document}
